\documentclass[11pt,a4paper]{article}
\usepackage[margin=25mm]{geometry}
\usepackage[T1]{fontenc}
\usepackage{hyperref}
\newcommand{\code}[1]{\texttt{\detokenize{#1}}}
\title{Blame: Open-Source Readiness}
\author{Blame Engineering}
\date{October 9, 2026}
\begin{document}
\maketitle

\section{Release direction and current status}
A public codebase can support a paid managed service. Users can run the software
with their own infrastructure and provider credentials, while the hosted offer
charges for operation, convenience, and metered services funded by the operator.
The application now separates free core access from Premium services funded by
Blame. This billing boundary does not yet make the project a complete,
documented self-hosting distribution.

A license decision is still required. A noncommercial restriction makes code
source-available rather than open source under the Open Source Definition.
A copyleft license such as AGPLv3 permits commercial use while imposing sharing
obligations in its applicable circumstances. The public license, contributor
terms, dependency compatibility, trademarks, and hosted-service terms must be
settled explicitly before release. No license choice is introduced by this pass.

\section{Separate the application from managed services}
Core projects, backlog, conversations, local worktrees, local previews, and
user-supplied AI credentials are available without a Blame subscription.
Authentication, authorization, and data ownership still apply. The global
subscription screen and core-route license middleware have been removed.

Premium applies to cloud voice transcription, managed Vercel previews, and new
hosted file uploads, including deployment source archives. Explicit service
checks protect those operations and cloud start or resume commands. Existing
inline attachment ingestion also checks Premium, while the administrator's
legacy-data migration can preserve previously stored attachments.

Existing file reads and attachment references remain available after expiry.
Cloud stop, suspend, deletion, and source recovery remain available so users can
clean up resources and recover their work. Connecting a personal Vercel account,
configuring its deployment target, and reading deployment metadata remain free;
new deployments through Blame require Premium for their hosted source upload.

The existing transcription and preview quotas, usage accounting, and Stripe
price identifiers are unchanged. The monthly offer is displayed as \$20;
annual pricing is confirmed by the existing Stripe checkout. An explicit
Premium page replaces the application-wide paywall, with upgrade actions on
paid operations. After expiry, active cloud preview controls remain available
for cleanup. A saved cloud preference falls back to local execution once no
cloud preview remains active, without deleting the saved choice.

A complete open-source distribution does not require hiding these checks in a
closed repository: the managed server can enforce its own policies even when
its implementation is public. A separate self-hosted configuration for
operator-funded services still needs an explicit design.

\section{Self-hosting dependencies to resolve}
\begin{itemize}
\item Make the backend origin and desktop authentication callback flow
configurable. Packaged clients currently assume the Blame cloud origin, with
development exceptions.
\item Replace the development command's personal SSH database tunnel with
documented local services and environment templates.
\item Document MySQL setup and migrations. The backend uses MySQL-specific
queries, JSON behavior, and locking; SQLite is not an interchangeable default.
\item Provide durable object storage configuration. The existing encrypted S3
adapter and ephemeral memory implementation do not constitute a durable local
filesystem adapter.
\item Make email delivery and verification setup explicit. Production
authentication currently requires email verification and an email provider.
\item Review required telemetry configuration so running a private instance
does not require the project's hosted telemetry account.
\item Expose the availability and configuration of speech transcription. The
current backend integration uses an operator OpenRouter credential; it is not
yet a local speech engine or a per-user transcription adapter.
\item Document cloud preview and Vercel deployment credentials separately.
Sandbox previews and deployment to a user's account are different workflows.
\end{itemize}

Local previews and user-supplied AI credentials already provide useful portable
capabilities. Deployment instructions should state which features run locally,
which require external credentials, and which are offered by the managed
service. They must not promise zero infrastructure or inference cost.

\section{Practical release sequence}
The global billing dependency has been removed from core flows and replaced
with service-specific checks. Next establish a reproducible local backend
setup, configurable origins, storage, email, and telemetry. Verify the default
clone-to-development workflow without personal infrastructure. Finally choose
the license, document contribution and security reporting, and validate desktop
packaging and provider setup across supported platforms.

The code refactoring and billing boundary are completed maintainability and
product steps in that sequence. Self-hosting configuration, public release,
and hosted-service publication remain separate implementation and release work.
\end{document}
